\begin{abstract}
Traditional weight of evidence assigns a credit characteristic a value based on its association with default across a population. Credit analysis also requires an assessment of that characteristic in the context of an individual applicant's record. We study conditional weight of evidence for this purpose. It measures the change in default log odds from observing a characteristic after the remaining credit information is known. Its sign identifies whether the characteristic adds adverse or favorable information for that applicant, and its magnitude measures the strength of that information. The same characteristic value can therefore receive different assessments for applicants with different payment histories or other credit attributes. We develop estimators using gradient-boosted trees and evaluate their accuracy in a bureau-style simulation. Regressions of the full model's predicted probabilities recover conditional evidence more accurately than separate models fitted to default outcomes. They also recover a deliberately constructed file-age effect whose direction varies with utilization and is often missed by marginal weight of evidence. Local class-frequency estimates improve on marginal weight of evidence but remain inaccurate for some weak effects. An application to FICO HELOC data illustrates estimation when the population evidence is unknown. Conditional weight of evidence provides an applicant-specific interpretation of the information in a credit characteristic. Its relation to adverse action depends additionally on how that information contributes to the fitted score difference between the applicant and the chosen passing records.
\end{abstract}
